\section{Boundary diagnostic and the unproved fixed-class comparison}

\begin{example}\label{ex:p143n-zero-class}
Assume \(\dim X=n\geq1\), and take \(\alpha=0\), \(\theta=0\).
Fix \(c\in\mathbf R\). At every K\"ahler perturbation \(\eta\),
let \(\Gamma^{\eta}_\tau=0\) for \(\tau<c\), with the usual
terminal and post-terminal extension. Since \(\eta\) is positive,
\(P_\eta[0]_{\mathcal I}=0\), and these curves are compatible.
All multiplier ideals are \(\mathcal O_X\), so the raw function is
the constant \(c\). The nonpluripolar mass of each component is
\(\int_X\eta^n>0\), but the mass of the inverse-system object is
\[
 \lim_{\varepsilon\downarrow0}
 \int_X(\varepsilon\omega)^n
 =\lim_{\varepsilon\downarrow0}
 \varepsilon^n\int_X\omega^n=0.
\]
There is no positive-mass unperturbed zero-class curve to which a
positive-volume recovery lemma could apply. Nevertheless the object
and its raw function are defined without ambiguity, and the constant
\(c\) belongs to every \(\PSH_{\mathrm{MP}}(\varepsilon\beta)\).
Thus zero volume of the limiting class does not obstruct the
inverse-system argument. This example does not classify all of
\(\mathcal Q_\omega(0)\).
\end{example}

For clarity, one can separately define an auxiliary fixed-class model
completion \(\mathcal C_{\mathrm{mod}}(\alpha)\), even for nef
boundary \(\alpha\): its members are nontrivial decreasing limits
of model functions \(\varphi_D\) for which
\(\alpha_{\mathcal X}+D\) is relatively nef on a determining
test configuration. The same-model argument used in
\eqref{eq:p143n-mp-monotone} proves only
\[
 \mathcal C_{\mathrm{mod}}(\alpha)
 \subset\mathcal Q_\omega(\alpha).
\]
The reverse inclusion would require approximants in the fixed class
\(\alpha\), not separate nets in every larger class. Our diagonal
construction produces compatible DXZ components and does not provide
that approximation. No equality or strictness of this last inclusion
is asserted here.

If \(\alpha\) is pseudoeffective but non-nef, not all arbitrarily
small K\"ahler perturbations \(\alpha+\varepsilon\beta\) are
K\"ahler. The K\"ahler comparison hypothesis therefore cannot be
applied on the cofinal ray used in the proof. Replacing that ray by
only its large K\"ahler perturbations is not cofinal toward zero
and does not repair the argument. This identifies the exact boundary
of this deduction, rather than asserting a general pseudoeffective
comparison or a zero-volume Monge--Amp\`ere or energy theory.
